\chapter{Chốt Tiết Có Dư Vị}
\chapterintro{Cuối giờ không chỉ là hết giờ}{Nếu chốt đúng, tiết học không rơi cái bịch xuống tiếng trống mà còn ở lại trong đầu học sinh.}

\section{Một câu chốt cả lớp nhớ}
\begin{tinhhuong}{Tình huống}
Tiết học xong nhưng nếu kết lại nhạt, học sinh sẽ đi ra ngoài với cảm giác rất rời.
\end{tinhhuong}
\begin{goiydungngay}
Kết bằng cấu trúc: \textit{``Nếu hôm nay chỉ giữ lại một câu, thầy cô muốn lớp giữ câu này...''}
\begin{itemize}
\item Câu chốt phải ngắn.
\item Đừng dùng hai ba ý cùng lúc.
\item Có thể cho cả lớp nhắc lại một lần.
\end{itemize}
\end{goiydungngay}
\begin{cauchot}
Chốt ít nhưng găm sâu tốt hơn chốt nhiều mà tan nhanh.
\end{cauchot}
\ngancach

\section{Phiếu ra cửa bằng miệng}
\begin{tinhhuong}{Tình huống}
Cuối giờ muốn kiểm tra nhanh xem lớp còn giữ được gì, nhưng không đủ thời gian phát giấy.
\end{tinhhuong}
\begin{goiydungngay}
Khi học sinh chuẩn bị ra, nói: \textit{``Ra cửa bằng một câu thôi: hôm nay em nhớ nhất gì?''}
\begin{itemize}
\item Cho 5 đến 7 em trả lời ngắn.
\item Không sửa lỗi dài dòng ở cuối tiết.
\item Ghi nhanh lại vài câu hay để dùng cho tiết sau.
\end{itemize}
\end{goiydungngay}
\begin{cauchot}
Cuối giờ không cần kiểm tra đủ, chỉ cần biết bài học còn ở lại chỗ nào.
\end{cauchot}
\ngancach

\section{Điều em mang về hôm nay}
\begin{tinhhuong}{Tình huống}
Muốn chốt tiết theo hướng nhẹ, có cảm xúc, không biến cuối giờ thành thêm một lần hỏi bài.
\end{tinhhuong}
\begin{goiydungngay}
Hỏi cả lớp: \textit{``Không cần đúng sách, chỉ cần đúng em: hôm nay em mang về điều gì?''}
\begin{itemize}
\item Có thể là một ý, một câu, một cảm giác.
\item Chấp nhận câu trả lời đời thường.
\item Đây là cách chốt rất hợp giờ chủ nhiệm hoặc cuối tuần.
\end{itemize}
\end{goiydungngay}
\begin{cauchot}
Bài học đi xa hơn khi học sinh thấy nó có gì để mang về.
\end{cauchot}
\ngancach

\section{Ba chữ trước tiếng trống}
\begin{tinhhuong}{Tình huống}
Cuối tiết không còn nhiều thời gian, nhưng thầy cô vẫn muốn lớp giữ lại một dấu mốc thật gọn.
\end{tinhhuong}
\begin{goiydungngay}
Nói: \textit{``Trước tiếng trống, mỗi em chọn đúng ba chữ để nhớ về tiết này.''}
\begin{itemize}
\item Có thể cho học sinh nói miệng hoặc ghi nhanh.
\item Ba chữ buộc các em phải cô lại ý chính.
\item Rất hợp với các lớp đông.
\end{itemize}
\end{goiydungngay}
\begin{cauchot}
Càng ít chữ, học sinh càng phải nhớ điều thật đáng giữ.
\end{cauchot}
\ngancach

\section{Lời hẹn một dòng cho ngày mai}
\begin{tinhhuong}{Tình huống}
Muốn cuối tiết không chỉ khép lại mà còn mở đường cho buổi học sau.
\end{tinhhuong}
\begin{goiydungngay}
Kết bằng một lời hẹn rõ: \textit{``Ngày mai, ai mang đến lớp một ví dụ gần nhất với bài hôm nay sẽ mở tiết giúp thầy cô.''}
\begin{itemize}
\item Lời hẹn phải nhỏ và làm được thật.
\item Đừng biến nó thành một bài tập lớn.
\item Tiết sau sẽ vào bài nhanh hơn nhiều.
\end{itemize}
\end{goiydungngay}
\begin{cauchot}
Cuối tiết đẹp là cuối tiết khiến ngày mai bắt đầu dễ hơn.
\end{cauchot}
